\section{Probing, nuisance and intervention details}
\label{app:probing}
\paragraph{Activation capture.} Activations are read at the anchor position by forward hooks on the output of each decoder block. Prompts are processed in batches, left-padded so that the anchor has the same index in every row. In the same forward pass we append the gold answer tokens after the anchor. Because attention is causal, this leaves the anchor's activations unaffected, but it gives the log-probability the model assigns to the gold answer, which we use as a nuisance covariate. The cached keys and values are then cut back to the prompt, and the model decodes greedily from there. Layer index $l$ denotes the residual stream after block $l$ ($l=0$: embeddings).

\paragraph{Probes.} Probes are \texttt{scikit-learn} logistic regressions with L2 regularisation ($C=0.5$) and balanced class weights, on features standardised on the training split \cite{alain2016probes}; the layer is selected by grouped 3-fold cross-validation on training questions (never on test), and the reported number is the test AUROC at that layer. The probe direction is the weight vector $w$ divided by the per-feature standard deviations $\sigma$, so that it lives in the original activation space, and normalised to unit norm. The difference-of-means direction is the normalised difference of class means of $\Delta h$ (or of $h$ for state tasks). Bootstrap intervals resample test questions (1{,}000 draws) using the stored probe weights.

\paragraph{State probes.} State probes use the stored anchor activation of each condition, one example per condition, with the same estimator, standardisation, layer selection and question-level splits as the update probes. For a two-hop question the state-sufficiency task has five positive conditions: the minimal sufficient context in each hop order (de-duplicated when identical), sufficient plus a distractor, sufficient plus a duplicate, and the full context. It has ten negative ones: no evidence, distractors only, each single hop, and each penultimate state plus a distractor, a duplicate or an answer-string control. Conditions containing a falsified paragraph are excluded. All positives except the minimal sufficient context contain three or more paragraphs, and no negative does, so the task is partly separable by paragraph count; \cref{app:behaviour} reports a count-matched version restricted to two-paragraph contexts. State-correctness and abstention-flag probes use every condition, labelled by the model's own answer.

\paragraph{Nuisance covariates.} $\Delta$tokens, total tokens, $\Delta$ gold log-probability, $\Delta$ next-token entropy, $\Delta$ top probability, answer string in added text / in the full context / in the previous context, TF-IDF cosine relevance of the added passage to the question, number and word count of added paragraphs, relative insertion position, and a 16-dimensional truncated-SVD embedding of the added text's TF-IDF vector (fit on training increments). The nuisance-only classifier is a logistic regression on these covariates. The residualised probe first regresses $\Delta h$ on the covariates with ridge regression ($\lambda=1$) and then probes what the covariates do not explain. The ``gain'' is the AUROC added when out-of-fold probe scores (5-fold, grouped by question) are given to the nuisance classifier as an extra feature.

\paragraph{Text baselines.} TF-IDF (1--2-grams, sublinear tf, $\le200$k features) with logistic regression; DeBERTa-v3-base fine-tuned with AdamW (lr $3{\times}10^{-5}$, linear decay, class-weighted cross-entropy). The reduced budget uses one epoch, 256 tokens and 8{,}000 training increments (all positives kept); the full budget two epochs, 512 tokens and all training increments.

\paragraph{Interventions.} Steering adds the vector $\alpha\,u\,v$ at the anchor and every generated position (``anchor onward'') or at all positions; the unit $u$ is $|\mathbb E[v^\top h\mid\text{sufficient}]-\mathbb E[v^\top h\mid\text{insufficient}]|$ on training conditions, computed per layer, so that $\alpha=1$ moves the activation by one gap between the two class means. Projection-out sets $v^\top h$ to the insufficient-state mean at the anchor and generated positions for every layer from the best layer onward. Patching replaces the anchor residual during the prefill at every layer from the best layer to the last; decoding then proceeds from the patched cache. Principal components are computed from 1{,}500 training decisive updates per layer; the ``unembedding'' direction is the normalised row of the output embedding matrix for the gold answer's first token.
